\section{Conclusion}
\label{sec:conclusion}

We set out to give an EEG model the one relation its tokens never carry,
the coupling between frequency bands, and to measure what that buys. Two
things did. A small model whose tokens are learned frequency bands and
whose encoder attends along electrodes, bands and time in turn is
competitive with foundation models twenty to forty times its size on twelve
clinical corpora and far ahead of them on seizure detection, and
architectural controls show that both ingredients are necessary. The
coupling token itself is decisive on one kind of label, epileptiform events
whose morphology is a slow rhythm carrying fast activity, where it adds
$0.15$ in Cohen's $\kappa$ inside our encoder and transfers into CBraMod's
with the gain attributable to the coupling content; on every other label we
tested it is neutral, and we report it as such. Pretraining, evaluated to
convergence with probes, label-fraction curves and clean-transfer
controls, helps this model where labels are scarce and constrains it where
they are not. The practical conclusion is that cross-frequency structure
is worth building into a clinical EEG model, and that where it is built in
--- as token content or as an axis of the encoder --- should follow the
label the model is asked to predict.
